\documentclass[11pt]{article}

\usepackage[a4paper,margin=2.2cm]{geometry}
\usepackage[T1]{fontenc}
\usepackage{lmodern}
\usepackage{enumitem}
\usepackage{xcolor}
\usepackage{titlesec}
\usepackage{listings}
\usepackage{fancyhdr}
\usepackage{hyperref}

\definecolor{accent}{HTML}{1D9BF0}
\definecolor{codebg}{HTML}{F5F7F9}
\definecolor{codekw}{HTML}{7A3E9D}

\hypersetup{colorlinks=true, linkcolor=accent, urlcolor=accent}

\lstset{
  basicstyle=\ttfamily\footnotesize,
  backgroundcolor=\color{codebg},
  keywordstyle=\color{codekw}\bfseries,
  breaklines=true,
  frame=single,
  rulecolor=\color{gray!30},
  columns=fullflexible,
  language=Java,
  showstringspaces=false,
  xleftmargin=6pt,
  xrightmargin=6pt,
  aboveskip=6pt,
  belowskip=4pt
}

\titleformat{\section}{\large\bfseries\color{accent}}{\thesection.}{0.6em}{}
\titlespacing*{\section}{0pt}{14pt}{6pt}

\pagestyle{fancy}
\fancyhf{}
\rhead{\small Twitter Clone --- Oral Examination}
\lhead{\small Advanced Programming, Summer 2026}
\cfoot{\thepage}
\renewcommand{\headrulewidth}{0.4pt}

\newcounter{q}
\newcommand{\question}[1]{\refstepcounter{q}\item[\textbf{Q\theq.}] #1}

\setlist[description]{leftmargin=2.4em, labelsep=0.5em, itemsep=7pt}

\begin{document}

\begin{center}
  {\LARGE\bfseries Twitter Clone --- Code Defense Questions}\\[4pt]
  {\large Oral examination / project handover}\\[2pt]
  {\color{gray}Advanced Programming --- Dr.\ Saeed Reza Kheradpisheh}
\end{center}

\vspace{4pt}
\noindent
The following 30 questions are designed to verify that the team understands
every layer of its own implementation: the socket layer and protocol, server
concurrency, the JDBC/DAO layer and SQL, object-oriented design, and the JavaFX
client. Each student should be able to justify the specific design choices in the
code, not merely describe what it does.

%======================================================================
\section{Architecture \& Networking}
\begin{description}

\question On the server, each connection is handled by a \texttt{ClientHandler}
that \texttt{implements Runnable} and is submitted to a fixed thread pool created
with \texttt{Executors.newFixedThreadPool(20)}. Why a bounded pool rather than
creating a new \texttt{Thread} per client, and what happens when the pool is
exhausted?

\question The protocol reads and writes one JSON object \emph{per line} using
\texttt{BufferedReader.readLine()} and \texttt{PrintWriter.println()}. What
assumption does this make about the message content, and how would the protocol
break if a tweet's text contained a newline character? How is that avoided?

\question Every request carries a \texttt{requestId} (a UUID) that the server
echoes back on the response. What concrete problem does this solve on the client
compared to matching responses by their packet \texttt{type}?

\question In \texttt{ClientHandler.handleLine}, the call to
\texttt{Dispatcher.dispatch} is wrapped in a \texttt{try/catch (RuntimeException)}.
Why catch a \texttt{RuntimeException} here specifically, and what would happen to
the client's connection if this catch were removed and a handler threw an
unchecked exception?

\question Real-time delivery is implemented by \texttt{ConnectionRegistry}, a
class with only \texttt{static} members and a private constructor. Walk through
what happens, step by step, from the moment user A posts a tweet to the moment it
appears in user B's feed. Which users receive the \texttt{NEW\_TWEET\_PUSH}?

\question A push message (e.g.\ \texttt{NEW\_TWEET\_PUSH}) has no
\texttt{requestId}, while a response does. On the client, how does
\texttt{NetworkService} use this difference to decide whether to invoke a pending
callback or a push listener?

\end{description}

%======================================================================
\section{Concurrency \& Thread Safety}
\begin{description}

\question \texttt{Authenticator} and \texttt{ConnectionRegistry} both store their
state in a \texttt{ConcurrentHashMap}. Which threads access these maps
simultaneously, and what specific bug could occur if a plain \texttt{HashMap}
were used instead?

\question \texttt{ClientHandler.sendPacket} is declared \texttt{synchronized}.
Given that each client has its own handler, which two events could otherwise write
to the same socket concurrently, and what would the symptom of that race be on the
wire?

\question On the client, the background listener thread never updates the UI
directly --- it wraps every callback in \texttt{Platform.runLater(...)}. What rule
of JavaFX does this respect, and what exception would you expect if the UI were
updated from the listener thread instead?

\question The database is accessed through a HikariCP connection \emph{pool}
exposed as a singleton \texttt{DatabaseConnection}. Why is a pool preferable to
opening a fresh \texttt{DriverManager.getConnection()} in each DAO method, and how
does \texttt{try-with-resources} interact with the pool (does \texttt{close()}
really close the socket to Postgres)?

\end{description}

%======================================================================
\section{Database, DAO \& SQL}
\begin{description}

\question Every DAO builds its queries with \texttt{PreparedStatement} and
parameter binding (\texttt{setInt}, \texttt{setString}) rather than string
concatenation. Show, with a concrete example from \texttt{SearchDAO}, how this
prevents SQL injection in the \texttt{ILIKE} search.

\question \texttt{TweetDAO.createTweet} sets \texttt{conn.setAutoCommit(false)},
performs three inserts (tweet, media, hashtags), then commits --- and rolls back
on any \texttt{SQLException}. Why must these three inserts be one transaction, and
what inconsistent state could arise without it?

\question The feed query flattens retweets using
\texttt{JOIN tweets d ON d.id = COALESCE(r.retweet\_of, r.id)}. Explain what
\texttt{r} and \texttt{d} represent, and why \texttt{COALESCE} produces the
correct "display tweet" for both an original tweet and a repost.

\question The tweet-reading queries compute \texttt{liked} and \texttt{retweeted}
per row with correlated \texttt{EXISTS} subqueries parameterized by the viewer's
id. Why compute these on the server in one query instead of returning raw tweets
and letting the client ask separately?

\question \texttt{FollowDAO.follow} and \texttt{LikeDAO.like} use
\texttt{INSERT ... ON CONFLICT DO NOTHING}. What property does this give the
operation, and why does it matter if the same "Follow" button is clicked twice
quickly by the same user?

\question The \texttt{follows} and \texttt{likes} tables use a \emph{composite}
primary key \texttt{(follower\_id, following\_id)} / \texttt{(user\_id, tweet\_id)}
instead of a surrogate \texttt{SERIAL id}. Justify this choice.

\question \texttt{TweetDAO.deleteTweet} uses a \texttt{WITH RECURSIVE} CTE to
collect a tweet and its entire reply subtree, then deletes their likes and the
tweets themselves. Why is the recursion necessary, and why can the parent and its
child rows be removed in a single \texttt{DELETE ... WHERE id = ANY(?)} without a
foreign-key violation?

\question The schema was extended \emph{additively} (using
\texttt{ADD COLUMN IF NOT EXISTS} and \texttt{CREATE TABLE IF NOT EXISTS}) and is
re-run on every server startup by \texttt{SchemaInitializer}. What does
"idempotent" mean here, and why is re-running the whole script on each boot safe?

\question \texttt{MediaDAO.getForTweets} and \texttt{HashtagDAO.getForTweets}
fetch data for many tweets in one query using \texttt{WHERE tweet\_id = ANY(?)}
with a \texttt{java.sql.Array}. What performance problem (hint: "N+1") does this
avoid compared to querying inside the row-mapping loop?

\question Passwords are stored via \texttt{PasswordUtil.hash} which calls
\texttt{BCrypt.hashpw(plain, BCrypt.gensalt())}. Where is the salt stored, why is
BCrypt preferable to a plain SHA-256 hash, and why can two users with the same
password have different \texttt{password\_hash} values?

\end{description}

%======================================================================
\section{Protocol \& Serialization}
\begin{description}

\question The \texttt{Packet.payload} field is typed as \texttt{JsonElement}, but
the getter returns a \texttt{JsonObject}. This was changed from a plain
\texttt{JsonObject} field. What runtime error did the original type cause when a
message arrived with \texttt{"payload": null}, and how does the current design fix
it?

\question The \texttt{User} model marks \texttt{passwordHash} as \texttt{transient}.
What does \texttt{transient} do during Gson serialization, and why is this a
security-relevant decision when a \texttt{User} is sent to the client?

\question \texttt{Packet} exposes static factory methods \texttt{request},
\texttt{ok}, \texttt{error}, and \texttt{push}. What is the benefit of these
factories over calling the constructor directly, and which one attaches a
\texttt{requestId}?

\question The tweet character limit (280) lives in
\texttt{shared/Protocol.MAX\_TWEET\_LENGTH} and is enforced both in the JavaFX
\texttt{Composer} and again in \texttt{TweetHandler.handleCreateTweet}. Why
validate on \emph{both} sides rather than trusting the client?

\end{description}

%======================================================================
\section{Object-Oriented Design \& Patterns}
\begin{description}

\question Name the design pattern behind each of the following and explain why it
fits: (a) \texttt{DatabaseConnection}/\texttt{NetworkService}/\texttt{UserContext};
(b) \texttt{UserDAO}/\texttt{TweetDAO}/\ldots; (c) \texttt{Dispatcher};
(d) \texttt{ConnectionRegistry} together with the client's push listeners.

\question The \texttt{Dispatcher} centralizes authentication: for every packet
type except \texttt{PING}/\texttt{REGISTER}/\texttt{LOGIN} it calls
\texttt{Authenticator.validate(token)} and passes the resolved \texttt{userId} to
the handler. What are the advantages of doing this in one place versus checking
the token inside each handler?

\question \texttt{TweetDAO} defines a small functional interface
\texttt{Binder} so that a single private \texttt{query()} method can serve every
read (feed, profile, replies, search) by letting each caller bind its own
parameters. Which OOP principle does this illustrate, and what would the code look
like without it?

\end{description}

%======================================================================
\section{JavaFX Client}
\begin{description}

\question A \texttt{TweetCard} performs an "optimistic" like: it increments the
count in the UI immediately, then sends \texttt{LIKE\_TWEET} and \emph{reverts} if
the server responds with an error. What is the UX benefit, and how does the card
reconcile its state when a \texttt{LIKE\_PUSH} later arrives from another user's
action?

\question Dark mode is implemented by adding/removing a \texttt{"dark"} style
class on the scene root (\texttt{Theme.refresh}), and the stylesheet defines CSS
variables like \texttt{-fx-bg} under \texttt{.root} and \texttt{.root.dark}. Why
is this approach preferable to setting inline styles on individual controls when
the theme toggles?

\question Image attachments are read from disk, Base64-encoded into a
\texttt{data:} URI in the \texttt{Composer}, sent inside the JSON payload, stored
as \texttt{TEXT} in the \texttt{media} table, and decoded back with
\texttt{UiUtils.decodeImage}. Discuss the trade-offs of embedding images as Base64
in the database versus storing files on disk and keeping only a path/URL. What is
the effect on payload size, and why does the code cap the encoded length?

\end{description}

\vspace{10pt}
\hrule
\vspace{6pt}
\noindent{\small\itshape Tip for graders: ask the student to open the referenced
file and point to the exact lines while answering. Strong answers cite the
\emph{reason} for a choice (safety, concurrency, performance, UX), not just the
mechanics.}

\end{document}
